\section{The rules dialect, clause by clause}
\label{app:dialect}

Six of the eight contested rule choices of \S\ref{sec:game-dialect} are engine switches, so their
effect is measurable rather than assumed; the misdeclaration award is hardcoded, and the
forced-endgame ladder has a threshold flag but not one the scored harness exposes, so within a
scored match it moves only inside the \texttt{-{}-legacy} bundle. This appendix lists each choice
against the switch that changes it, so that a reader can re-implement the dialect or vary it.

\begin{center}\tabfont
\begin{tabular}{llp{0.42\linewidth}}
\toprule
Clause & Switch & Setting studied \\
\midrule
Deck and half-suits & \texttt{-{}-sets=8} & 54 cards in nine half-suits; \texttt{-{}-sets=8} plays the
48-card, eight-half-suit variant \\
Out-of-turn declaration & \texttt{-{}-no-out-of-turn} & permitted; the switch forbids it \\
Cardless declaration & \texttt{-{}-no-cardless-declare} & permitted; the switch forbids it \\
Simultaneous declarers & \texttt{-{}-arb=high}, \texttt{-{}-arb=turn} & lowest seat index first;
the switches arbitrate by stated confidence and by turn order \\
Misdeclaration & --- & the half-suit is awarded to the opposing team \\
Forced endgame & --- & eight-level willingness ladder when one whole team is cardless \\
Action cap & \texttt{-{}-maxasks=N} & 400 asks; residue adjudicated by physical majority \\
Legacy bundle & \texttt{-{}-legacy} & the full v0.3-era dialect, bundling several of the above \\
\bottomrule
\end{tabular}
\end{center}

\S\ref{sec:results-dialects} measures the paper's headline comparison over \vsevenDiaRows{}
dialect rows: the default, and seven variations drawn from six of these clauses. Every row keeps its sign and every row replicates; the largest excursion from the
default row is \vsevenDiaMaxExcursion{}~pp, on \vsevenDiaMaxExcursionRow{}.

\paragraph{Two clauses on which an earlier report in this corpus is wrong, corrected here.}

\begin{enumerate}
\item \textbf{An undeclared half-suit does not score nothing.} An earlier report justified an
  unconditional forcing rule on the premise that an unclaimed half-suit is worthless. Under the
  dialect actually evaluated, residue at the action cap is awarded by physical majority, and on a
  wrong post-horizon declaration the declaring team holds a mean \numResidueHeld{} of the six cards.
  The rule converts a probable award into a certain concession.
\item \textbf{Restarting the forced-endgame ladder sharpens nothing.} The argument is that restarting
  the willingness ladder after each declaration lets early, safer declarations resolve cards and so
  sharpen later ones. It is sound and, in the play it was measured on, inoperative: on every occasion
  a team went cardless with live half-suits --- \numLadderOccasions{} such records at seed
  \numLadderSeed{} in \vfour{} mirror self-play, half of them duplicates, because at each deal shift
  the two team orientations play identical policies and so record every event twice ---
  exactly \numLadderLiveSets{} half-suit was live,
  so there was no ordering to exploit. That is a strong tendency and not a universal: a second
  sample at seed \vsevenLadderSecondSeed{} finds \vsevenLadderSecondTotal{} \emph{distinct} forced
  declarations, all but one of them with a single live half-suit. It is also
  policy-dependent: the source
  records that a policy made more patient about declaring would begin to produce multi-half-suit
  forced endgames, at which point the ordering would start to matter. \fishseven{} is a change in
  that direction, since it disables the urgency escalation entirely
  (\S\ref{sec:agent-urgency}), and no measurement of live-half-suit multiplicity at forced-endgame
  entry exists for it.
\end{enumerate}
